\documentclass[11pt,a4paper]{article}

\usepackage[margin=2.5cm]{geometry}
\usepackage{amsmath,amssymb,amsthm}
\usepackage{booktabs,array,multirow}
\usepackage{enumitem}
\usepackage{xcolor}
\usepackage{algorithm}
\usepackage{algpseudocode}
\usepackage[colorlinks=true,linkcolor=blue!60!black,citecolor=blue!60!black,urlcolor=blue!60!black]{hyperref}

\newtheorem{lemma}{Lemma}
\newtheorem{proposition}{Proposition}
\theoremstyle{remark}
\newtheorem{remark}{Remark}
\theoremstyle{definition}
\newtheorem{issue}{Ambiguity}

% ---------- notation macros ----------
\newcommand{\sem}{{261S}}   % semiconductors
\newcommand{\oth}{{26X}}    % ICT & electronics excl. semiconductors
\newcommand{\agg}{{26}}     % original division 26
\newcommand{\tZ}{\tilde{Z}}
\newcommand{\todo}[1]{\textcolor{red}{[\textbf{TODO:} #1]}}
\newcommand{\paper}[1]{\textcolor{blue!50!black}{[paper: #1]}}

\title{Splitting ICIO division 26 into semiconductors and other ICT \& electronics:\\
a step-by-step guide to Annex~A of Haramboure et al.\ (2023)}
\author{Working notes}
\date{\today}

\begin{document}
\maketitle

\begin{abstract}
\noindent
Haramboure, Lalanne, Schwellnus and Guilhoto (2023), \emph{Vulnerabilities in the semiconductor supply chain},
OECD STI Working Paper 2023/05, split ICIO industry ``Computer, electronic and optical equipment'' (ISIC Rev.4
division 26) into ``Semiconductors'' (\sem) and ``ICT and electronics excluding semiconductors'' (\oth). Their
Annex~A describes the method in four pages. This note restates every step with explicit notation, spells out the
algebra behind each equation, shows which data enter where, gives a worked numerical example of the balancing
step, and lists the points where the Annex is ambiguous or inconsistent together with the resolution adopted
here. Equation labels (A.1)--(A.13) refer to the transcription in \texttt{doc/haramboure2023\_annexA.md}.
\end{abstract}

\tableofcontents
\newpage

%==========================================================================
\section{What the method does, in one page}
%==========================================================================

\paragraph{Goal.} Start from a standard ICIO table in which division 26 is a single industry in every country.
End with a table in which, in every country, row and column 26 are replaced by two rows and two columns
(\sem{} and \oth), such that the new table is a complete, internally consistent ICIO table.

\paragraph{Key idea.} The paper does \emph{not} split flows directly. It splits the \emph{technical
coefficients} $A$, sets final demand for semiconductors to zero, and then recomputes gross output, flows and
value added from the Leontief model. Everything therefore hinges on constructing a split $A$ matrix.

\paragraph{Data.}
\begin{enumerate}[nosep]
  \item ICIO table (all years): $Z$, $Y$, $X$, value added, taxes.
  \item Bilateral trade in semiconductors (HS 8541 + 8542), all years and country pairs: $M^{r,s}_\sem$.
  \item A small number of detailed national IO tables (one base year each) in which semiconductors are a
        separate product/industry. The paper uses KOR 2014, JPN 2011, MYS 2010, TWN 2016, USA 2012.
\end{enumerate}

\paragraph{Pipeline.} The split $A$ matrix is built block by block (Figure~\ref{fig:blocks}):

\begin{center}
\begin{tabular}{@{}clll@{}}
\toprule
Step & Block of $A$ & Source of the split & Annex eq. \\
\midrule
0 & -- & Prepare ICIO, trade data, national-table ratios & -- \\
1 & domestic ($r=s$), reference economies & national IO tables & (A.1) \\
2 & domestic ($r=s$), other economies & mapping to reference economies & -- \\
3 & imports ($r\neq s$), columns & domestic column ratios of the \emph{buyer} $s$ & (A.2) \\
4 & imports ($r\neq s$), rows & domestic row shares of $s$ + RAS to match trade & (A.3)--(A.10) \\
5 & final demand, $X$, $Z$, VA & Leontief model & (A.11)--(A.13) \\
\bottomrule
\end{tabular}
\end{center}

\begin{figure}[ht]
\centering
\renewcommand{\arraystretch}{1.3}
\begin{tabular}{c|ccc|ccc|}
 \multicolumn{1}{c}{} & \multicolumn{3}{c}{buyer economy $s=r$} & \multicolumn{3}{c}{buyer economy $s\neq r$} \\
 \multicolumn{1}{c}{} & $j\neq 26$ & \sem & \multicolumn{1}{c}{\oth} & $j\neq 26$ & \sem & \multicolumn{1}{c}{\oth} \\ \cline{2-7}
 origin $r$, $i\neq 26$ & known & \multicolumn{2}{c|}{\textbf{col split} (1--2)} & known & \multicolumn{2}{c|}{\textbf{col split} (3)} \\
 origin $r$, \sem & \multirow{2}{*}{\textbf{row split} (1--2)} & \multicolumn{2}{c|}{\multirow{2}{*}{\textbf{2$\times$2} (1--2)}} & \multirow{2}{*}{\textbf{row split} (4)} & \multicolumn{2}{c|}{\multirow{2}{*}{\textbf{2$\times$2} (3+4)}}\\
 origin $r$, \oth & & & & & & \\ \cline{2-7}
\end{tabular}
\caption{Which cells of $A$ must be created, and by which step (numbers in brackets). ``Row split'': the
original row 26 is divided between \sem{} and \oth{} (who \emph{sells} what). ``Col split'': the original
column 26 is divided (who \emph{buys} inputs). The intersection is a $2\times 2$ block. This replaces the
paper's Figure A A.1, which is an image.}
\label{fig:blocks}
\end{figure}

%==========================================================================
\section{Notation and ICIO accounting}
%==========================================================================

\subsection{Indices and objects}

\begin{itemize}[nosep]
  \item Economies $r,s \in \{1,\dots,N\}$; $r$ = origin (seller), $s$ = destination (buyer).
  \item Industries $i,j \in \{1,\dots,K\}$ in the original table; $i$ = supplying (row), $j$ = purchasing (column).
        Division 26 is written \agg. After the split, \agg{} is replaced by \sem{} and \oth, so there are $K+1$ industries.
  \item $Z^{r,s}_{i,j}$: intermediate sales of industry $i$ in $r$ to industry $j$ in $s$ (USD million).
  \item $Y^{r,s}_{i}$: final demand of $s$ for output of industry $i$ in $r$, summed over final-demand
        categories (HFCE, NPISH, GGFC, GFCF, INVNT, DPABR).
  \item $X^r_i$: gross output. $W^r_j$: value added; $T^r_j$: taxes less subsidies on intermediates
        (the ICIO \texttt{TLS} row). Some ICIO editions record TLS by origin; treat it as one row here.
  \item Technical coefficients: $a^{r,s}_{i,j} = Z^{r,s}_{i,j} / X^s_j$.
        \emph{The denominator is the output of the buying industry.}
  \item Base year of a national table: $t_0$. A left superscript ${}^{t}a$ marks the year when needed.
  \item $M^{r,s}_\sem$: imports of semiconductors by $s$ from $r$ (trade data).
\end{itemize}

\subsection{Identities that the final table must satisfy}
\begin{align}
  \text{row (sales):}\quad & X^r_i = \sum_{s}\sum_j Z^{r,s}_{i,j} + \sum_s Y^{r,s}_i, \label{eq:row}\\
  \text{column (costs):}\quad & X^s_j = \sum_r\sum_i Z^{r,s}_{i,j} + T^s_j + W^s_j. \label{eq:col}
\end{align}
In matrix form with $A$ stacked over all $(r,i)$ and $(s,j)$: $X = AX + Y\mathbf 1$, hence $X = (I-A)^{-1}y$ with
$y$ the row sums of final demand.

\subsection{Two facts that explain every formula in Annex A}

\begin{lemma}[Rows add up]\label{lem:rows}
For any buyer $(s,j)$ and origin $r$,
$a^{r,s}_{\sem,j} + a^{r,s}_{\oth,j} = a^{r,s}_{\agg,j}$.
\end{lemma}
\begin{proof}
All three coefficients share the same denominator $X^s_j$, and $Z^{r,s}_{\sem,j}+Z^{r,s}_{\oth,j}=Z^{r,s}_{\agg,j}$.
\end{proof}

So splitting a \textbf{row} is splitting a number into two \emph{shares} that sum to one:
$a^{r,s}_{\sem,j} = \sigma\, a^{r,s}_{\agg,j}$, $a^{r,s}_{\oth,j} = (1-\sigma)\, a^{r,s}_{\agg,j}$.

\begin{lemma}[Columns average]\label{lem:cols}
Let $w^s = X^s_\sem / X^s_\agg$ be the semiconductor share of division-26 output in $s$. For any supplier $(r,i)$,
\[
  a^{r,s}_{i,\agg} = w^s\, a^{r,s}_{i,\sem} + (1-w^s)\, a^{r,s}_{i,\oth}.
\]
\end{lemma}
\begin{proof}
$a^{r,s}_{i,\agg} = \dfrac{Z^{r,s}_{i,\sem}+Z^{r,s}_{i,\oth}}{X^s_\sem + X^s_\oth}
 = \dfrac{X^s_\sem}{X^s_\agg}\dfrac{Z^{r,s}_{i,\sem}}{X^s_\sem} + \dfrac{X^s_\oth}{X^s_\agg}\dfrac{Z^{r,s}_{i,\oth}}{X^s_\oth}.$
\end{proof}

So splitting a \textbf{column} is \emph{not} a share split. The two new coefficients are not fractions of the old
one; they are \emph{intensities} whose output-weighted average is the old one. This is why Annex A uses ratios
such as $a_{i,\sem}/a_{i,\agg}$ (which can exceed 1) for columns and shares $a_{\sem,j}/a_{\agg,j}$ (which lie in
$[0,1]$) for rows.

\begin{remark}\label{rem:rho}
Define the column \emph{relative intensity} $\rho_{i,k} = a_{i,k}/a_{i,\agg}$ for $k\in\{\sem,\oth\}$.
Lemma~\ref{lem:cols} says $w\rho_{i,\sem} + (1-w)\rho_{i,\oth} = 1$. Equivalently, in flow terms,
the \emph{column flow share} $\phi_i = Z_{i,\sem}/Z_{i,\agg} = w\,\rho_{i,\sem}$ lies in $[0,1]$.
\end{remark}

%==========================================================================
\section{Step 0: preparing the inputs}
%==========================================================================

\subsection{ICIO}
Load $Z$, $Y$ (all final-demand categories), $X$, $W$, $T$ for every year. Check identities
\eqref{eq:row}--\eqref{eq:col} on the original data first; the tolerance you observe there is the tolerance you
can expect after the split. Compute $A$ with $a^{r,s}_{i,j} = Z^{r,s}_{i,j}/X^s_j$ (set to 0 when $X^s_j=0$).
If the ICIO edition splits China and Mexico (CN1/CN2, MX1/MX2), treat each part as an economy.

\subsection{Trade data: $M^{r,s}_\sem$}
\paper{BTDIxE database (\url{http://oe.cd/btd}), HS headings 8541 and 8542, every year and economy pair.}
\begin{itemize}[nosep]
  \item Units: convert to USD million to match ICIO.
  \item Valuation: trade data are FOB (exports) or CIF (imports); ICIO imports are at basic prices. The paper does
        not adjust. Decide and document (e.g.\ use exporter-reported FOB values, or mirror data).
  \item HS 8541 includes photovoltaic cells and LEDs. The paper keeps the whole headings in Annex A
        (it excludes PV and LEDs only in the descriptive Annex B figures).
  \item Aggregate trade partners to the ICIO country list (rest-of-world included).
  \item Only $r\neq s$ is used. Domestic sales of semiconductors are not observed.
\end{itemize}

\subsection{National IO tables: the three objects you need}
For each reference economy $c$ with a detailed table in base year $t_0$:
\begin{enumerate}
  \item Map the detailed products/industries to the ICIO industry list, but keep \sem{} and \oth{} separate
        (\sem{} = semiconductor devices / integrated circuits; \oth{} = the rest of division 26).
  \item Use the \textbf{domestic} flow table (intermediate use of domestically produced goods), because the
        paper uses domestic technical coefficients. If the table is of the competitive-import type (imports
        included in each cell), remove imports first, e.g.\ with the import-proportionality assumption.
        \todo{check table type for each national table you hold.}
  \item Compute base-year domestic coefficients ${}^{t_0}a^{c,c}_{i,j}$ on the $(K+1)\times(K+1)$ classification and
        on the original $K\times K$ classification (aggregate \sem+\oth{} back into \agg).
\end{enumerate}
From these, store:
\begin{align}
  \text{row shares:}\qquad & \sigma^c_j = \frac{{}^{t_0}a^{c,c}_{\sem,j}}{{}^{t_0}a^{c,c}_{\agg,j}}
     \quad\text{for all buyers } j \text{ (original } j \text{ and } j\in\{\sem,\oth\}), \label{eq:sigma}\\
  \text{column ratios:}\qquad & \rho^c_{i,k} = \frac{{}^{t_0}a^{c,c}_{i,k}}{{}^{t_0}a^{c,c}_{i,\agg}},
     \quad k\in\{\sem,\oth\}, \text{ for all suppliers } i, \label{eq:rho}\\
  \text{output share:}\qquad & w^c_0 = {}^{t_0}X^c_\sem / {}^{t_0}X^c_\agg . \label{eq:w0}
\end{align}
Also store the $2\times2$ domestic block ratios
$\beta^c_{k,l} = {}^{t_0}a^{c,c}_{k,l}/{}^{t_0}a^{c,c}_{\agg,\agg}$, $k,l\in\{\sem,\oth\}$,
and the value-added and tax ratios of the two new columns (needed in Step 5).
These ratios are unit-free, so national currency and differing base years are not a problem.

%==========================================================================
\section{Step 1: domestic blocks for reference economies, eq.~(A.1)}
%==========================================================================

\paper{``For other years, it is assumed that the shares of the base year can be used to split the technical
coefficients'' of division 26.}

For a reference economy $r$ and any year $t$, take the ICIO coefficients ${}^{t}a^{r,r}$ and apply the base-year
ratios:
\begin{align}
 \text{row split:}\quad & {}^{t}a^{r,r}_{\sem,j} = {}^{t}a^{r,r}_{\agg,j}\,\sigma^r_j, \qquad
                          {}^{t}a^{r,r}_{\oth,j} = {}^{t}a^{r,r}_{\agg,j}\,(1-\sigma^r_j), \tag{A.1a}\\
 \text{column split:}\quad & {}^{t}a^{r,r}_{i,k} = {}^{t}a^{r,r}_{i,\agg}\,\rho^r_{i,k},\qquad k\in\{\sem,\oth\}, \tag{A.1b}\\
 2\times 2\text{ block:}\quad & {}^{t}a^{r,r}_{k,l} = {}^{t}a^{r,r}_{\agg,\agg}\,\beta^r_{k,l},\qquad k,l\in\{\sem,\oth\}. \tag{A.1c}
\end{align}

\paragraph{Derivation / meaning.}
(A.1a) holds the base-year \emph{composition} of division-26 purchases by industry $j$ fixed: if 30\% of what
Korean motor vehicles bought from domestic division 26 in 2014 were chips, 30\% is assumed in every year.
(A.1b) holds the base-year \emph{relative input intensity} of the two sub-industries fixed: if the chip industry
used 1.8 times as much domestic electricity per unit of output as division 26 on average, that factor is kept.
(A.1c) follows from applying (A.1a) and (A.1b) in sequence:
${}^{t}a_{k,l} = {}^{t}a_{\agg,\agg}\,\rho_{\agg,l}\,\sigma^{(k)}_l
 = {}^{t}a_{\agg,\agg}\frac{{}^{t_0}a_{\agg,l}}{{}^{t_0}a_{\agg,\agg}}\frac{{}^{t_0}a_{k,l}}{{}^{t_0}a_{\agg,l}}
 = {}^{t}a_{\agg,\agg}\,\beta_{k,l}.$

\paragraph{Consistency.}
Rows aggregate exactly (Lemma~\ref{lem:rows}). Columns aggregate exactly \emph{with the base-year weight} $w^r_0$:
$w^r_0\rho^r_{i,\sem} + (1-w^r_0)\rho^r_{i,\oth} = 1$ by Remark~\ref{rem:rho}. In year $t$ the model will produce its
own $w^r_t$ (Step 5); if $w^r_t\neq w^r_0$ the split columns do not average back exactly to the ICIO column.
This is the main source of discrepancies checked in Section~\ref{sec:checks}.

\paragraph{Implementation.} In the base year itself (A.1) reproduces the national table's ratios but applies them
to ICIO levels, not national levels. You never copy national coefficient \emph{levels} into the ICIO.

%==========================================================================
\section{Step 2: domestic blocks for the other economies}
%==========================================================================

\paper{China $\leftarrow$ Malaysia; Singapore $\leftarrow$ mean(KOR, TWN); EU members $\leftarrow$ mean(JPN, USA);
all others $\leftarrow$ mean(JPN, KOR, MYS, TWN).}

For a non-reference economy $s$ with mapping set $\mathcal{R}(s)$, build ratios as averages over $c\in\mathcal{R}(s)$
and apply (A.1a)--(A.1c) to $s$'s own ICIO coefficients.

\begin{issue}[What is averaged]\label{iss:avg}
The paper says domestic technical \emph{coefficients} are a ``linear combination'' of those in the reference
economies. Averaging coefficient \emph{levels} would discard $s$'s own ICIO data. \textbf{Resolution:} average the
\emph{ratios}, and apply them to $s$'s own coefficients.
\end{issue}

\begin{issue}[Averaging column ratios breaks Lemma 2]\label{iss:colavg}
Row shares can be averaged freely: $\bar\sigma_j = \frac1{|\mathcal R|}\sum_c\sigma^c_j\in[0,1]$.
Column ratios cannot. Each $\rho^c$ satisfies $w^c_0\rho^c_{i,\sem}+(1-w^c_0)\rho^c_{i,\oth}=1$ with its own
$w^c_0$, so the simple average $\bar\rho$ satisfies no such identity with a single $w$.
\textbf{Resolution:} average the column \emph{flow} shares $\phi^c_i = w^c_0\rho^c_{i,\sem}$ and the output shares,
$\bar\phi_i = \mathrm{mean}_c\,\phi^c_i$, $\bar w = \mathrm{mean}_c\,w^c_0$, then set
\[
  \bar\rho_{i,\sem} = \bar\phi_i/\bar w,\qquad \bar\rho_{i,\oth} = (1-\bar\phi_i)/(1-\bar w),
\]
which satisfies $\bar w\bar\rho_{i,\sem}+(1-\bar w)\bar\rho_{i,\oth}=1$ exactly. Do the same for $\beta$.
\end{issue}

\begin{issue}[No Malaysian table]\label{iss:mys}
Malaysia enters the rules for China and for ``all others''. Options, in order of preference:
(a) obtain a Malaysian table with a semiconductor product (DOSM publishes IO tables; \todo{check whether a
semiconductor product is separately identified});
(b) use a Chinese national table if a sufficiently narrow ``electronic components'' sector exists
(\todo{check; this is broader than semiconductors});
(c) replace Malaysia by a proxy with a similar, assembly/test-heavy profile, e.g.\ China $\leftarrow$ mean(KOR, TWN)
and ``others'' $\leftarrow$ mean(JPN, KOR, TWN).
Report sensitivity of the results to this choice.
\end{issue}

After Steps 1--2, every diagonal block $A^{r,r}$ is split.

%==========================================================================
\section{Step 3: imported inputs, column split, eq.~(A.2)}
%==========================================================================

\paper{``Imported intermediate inputs of product $i$ in economy $s$ are split between the two sub-industries using
the shares observed for domestic inputs of product $i$ in economy $s$.''}

For every $r\neq s$ and every supplying industry $i$ of $r$:
\begin{equation}
  a^{r,s}_{i,\sem} = a^{r,s}_{i,\agg}\,\frac{a^{s,s}_{i,\sem}}{a^{s,s}_{i,\agg}} = a^{r,s}_{i,\agg}\,\rho^s_{i,\sem},
  \qquad
  a^{r,s}_{i,\oth} = a^{r,s}_{i,\agg}\,\frac{a^{s,s}_{i,\oth}}{a^{s,s}_{i,\agg}} = a^{r,s}_{i,\agg}\,\rho^s_{i,\oth}.
  \tag{A.2}
\end{equation}

\begin{issue}[Typo in the paper]
The paper prints $a^{r,s}_{i,\oth} = a^{r,s}_{i,\oth}\cdot(\dots)$, which is circular. The first factor must be
$a^{r,s}_{i,\agg}$, as in the \sem{} equation.
\end{issue}

\paragraph{Meaning.} The ratio of the \emph{buyer} $s$ is used, not of the seller $r$: how the Korean chip industry
uses a steel input relative to the rest of Korean division 26 is assumed to be the same whether the steel is
Korean or Japanese. This is the standard import-proportionality assumption applied to columns.

\paragraph{Flow-form equivalence.} Multiplying by $X^s_\sem = w^s X^s_\agg$:
\[
  Z^{r,s}_{i,\sem} = a^{r,s}_{i,\agg}\,\rho^s_{i,\sem}\,w^s X^s_\agg
   = Z^{r,s}_{i,\agg}\cdot \underbrace{w^s\rho^s_{i,\sem}}_{\phi^s_i}
   = Z^{r,s}_{i,\agg}\,\frac{Z^{s,s}_{i,\sem}}{Z^{s,s}_{i,\agg}}.
\]
So (A.2) is equivalent to ``split the imported flow in the same proportion as the domestic flow''. The coefficient
form has the advantage that it does not require $w^s$, which is unknown until Step 5.

\paragraph{Which rows $i$?} All supplying rows of $r$, \emph{including} the split rows \sem{} and \oth{} of $r$
once Step 4 has produced them. For those two rows use the $2\times2$ ratios of the buyer,
$a^{s,s}_{k,l}/a^{s,s}_{k,\agg}$, $k,l\in\{\sem,\oth\}$. This is where the order of Steps 3 and 4 matters
(Section~\ref{sec:order}).

%==========================================================================
\section{Step 4: imported division-26 goods, row split, eqs.~(A.3)--(A.10)}
%==========================================================================

This is the only step that uses trade data. It is carried out \textbf{separately for every ordered pair}
$(r,s)$, $r\neq s$, and every year. For a fixed pair, the object is a $2\times J$ matrix: two rows (\sem, \oth)
and one column per buying industry $j$ of $s$.

\subsection{Stage 1: initial guess, eq.~(A.3)}
Assume the semiconductor share of imported division-26 inputs equals the domestic one:
\begin{equation}
  Z^{r,s(0)}_{\sem,j} = \frac{a^{s,s}_{\sem,j}}{a^{s,s}_{\agg,j}}\,Z^{r,s}_{\agg,j} = \sigma^s_j\,Z^{r,s}_{\agg,j},
  \qquad
  Z^{r,s(0)}_{\oth,j} = Z^{r,s}_{\agg,j} - Z^{r,s(0)}_{\sem,j}. \tag{A.3}
\end{equation}
$Z^{r,s}_{\agg,j}$ is the observed ICIO flow; $\sigma^s_j$ comes from Steps 1--2. The paper works with flows here
(not coefficients) because the constraint to be matched, bilateral trade, is a flow.

\subsection{Stage 2: what must be balanced, eqs.~(A.4)--(A.5)}
Nothing guarantees $\sum_j Z^{r,s(0)}_{\sem,j} = M^{r,s}_\sem$ (A.4). The balanced matrix $\tZ$ must satisfy
\begin{align}
  \text{row 1 (trade):}\quad & \textstyle\sum_j \tZ^{r,s}_{\sem,j} = M^{r,s}_\sem, \label{eq:rowS}\\
  \text{row 2 (implied):}\quad & \textstyle\sum_j \tZ^{r,s}_{\oth,j} = \sum_j Z^{r,s}_{\agg,j} - M^{r,s}_\sem \equiv M^{r,s}_\oth, \label{eq:rowX}\\
  \text{columns (ICIO):}\quad & \tZ^{r,s}_{\sem,j} + \tZ^{r,s}_{\oth,j} = Z^{r,s}_{\agg,j}\quad\forall j. \label{eq:colj}
\end{align}

\begin{issue}[The \oth{} row target]
The paper states only \eqref{eq:rowS} and \eqref{eq:colj}, and says ``the same procedure'' is used for the \oth{}
row. The \oth{} row needs a target; the only one consistent with \eqref{eq:rowS}--\eqref{eq:colj} is
\eqref{eq:rowX}, obtained by summing \eqref{eq:colj} over $j$.
\end{issue}

\paragraph{Feasibility.} A non-negative solution exists only if $0\le M^{r,s}_\sem\le\sum_j Z^{r,s}_{\agg,j}$.
Since final demand for semiconductors is set to zero (Step 5), \emph{all} semiconductor imports must be absorbed by
intermediate division-26 imports. Trade data can exceed that total (valuation differences, re-exports, ROW
aggregation). \textbf{Rule:} cap $M^{r,s}_\sem$ at, e.g., $\sum_jZ^{r,s}_{\agg,j}$ minus a small margin, and log
every capped pair. Also, if $\sum_j Z^{r,s(0)}_{\sem,j}=0$ but $M^{r,s}_\sem>0$, RAS cannot scale up from zero;
seed with $\sigma^s_j$ replaced by the average of positive shares.

\subsection{Stage 2: the RAS iterations, eqs.~(A.6)--(A.9)}
The paper writes each RAS step as ``compute the error, allocate it proportionally''. With iteration counter $k$
(the paper uses $i$, which clashes with the industry index):

\emph{Row step.} Error $\varepsilon^{(k)}_\sem = M_\sem - \sum_j Z^{(k)}_{\sem,j}$ (A.6), and
\[
  Z^{(k+1)}_{\sem,j} = Z^{(k)}_{\sem,j} + \frac{Z^{(k)}_{\sem,j}}{\sum_{j'}Z^{(k)}_{\sem,j'}}\,\varepsilon^{(k)}_\sem
  \tag{A.7}
\]
(and the same for \oth{} with target $M_\oth$). Substituting the error:
\[
  Z^{(k+1)}_{\sem,j} = Z^{(k)}_{\sem,j}\left(1 + \frac{M_\sem - \sum_{j'}Z^{(k)}_{\sem,j'}}{\sum_{j'}Z^{(k)}_{\sem,j'}}\right)
  = Z^{(k)}_{\sem,j}\,\underbrace{\frac{M_\sem}{\sum_{j'}Z^{(k)}_{\sem,j'}}}_{r_\sem}.
\]
So the ``additive'' form is exactly a \textbf{multiplicative row scaling} by $r_\sem$: the classic $R$ of RAS.

\emph{Column step.} Error $\varepsilon^{(k+1)}_j = Z_{\agg,j} - Z^{(k+1)}_{\sem,j} - Z^{(k+1)}_{\oth,j}$ (A.8), and
\[
  Z^{(k+2)}_{\sem,j} = Z^{(k+1)}_{\sem,j} + \frac{Z^{(k+1)}_{\sem,j}}{Z^{(k+1)}_{\sem,j}+Z^{(k+1)}_{\oth,j}}\,\varepsilon^{(k+1)}_j
  = Z^{(k+1)}_{\sem,j}\,\underbrace{\frac{Z_{\agg,j}}{Z^{(k+1)}_{\sem,j}+Z^{(k+1)}_{\oth,j}}}_{s_j}.
  \tag{A.9}
\]
The same factor $s_j$ multiplies the \oth{} cell: the $S$ of RAS. After the column step the column constraints hold
exactly; after the row step the row constraints hold exactly. Alternate until both hold within tolerance.
Bacharach (1970) proves convergence for non-negative matrices when a solution exists. The limit is the matrix
closest to $Z^{(0)}$ in the relative-entropy sense that satisfies the constraints, so it preserves the
\emph{pattern} of the initial guess (zeros stay zero).

\subsection{Worked example}
One pair $(r,s)$, three buying industries $A,B,C$ with imported division-26 flows $Z_{\agg,\cdot} = (10,30,60)$,
total 100. Domestic shares in $s$: $\sigma = (0.1, 0.5, 0.2)$. Trade data: $M_\sem = 35$, hence $M_\oth = 65$.

\begin{center}
\begin{tabular}{@{}lrrrrr@{}}
\toprule
 & $A$ & $B$ & $C$ & row sum & target \\
\midrule
\multicolumn{6}{@{}l}{\emph{Stage 1, (A.3)}}\\
\sem & 1.000 & 15.000 & 12.000 & 28.000 & 35 \\
\oth & 9.000 & 15.000 & 48.000 & 72.000 & 65 \\
\midrule
\multicolumn{6}{@{}l}{\emph{Row step: $r_\sem = 35/28 = 1.25$, $r_\oth = 65/72 = 0.9028$}}\\
\sem & 1.250 & 18.750 & 15.000 & 35.000 & 35 \\
\oth & 8.125 & 13.542 & 43.333 & 65.000 & 65 \\
col sum & 9.375 & 32.292 & 58.333 & & \\
col target & 10 & 30 & 60 & & \\
\midrule
\multicolumn{6}{@{}l}{\emph{Column step: $s = (1.0667, 0.9290, 1.0286)$}}\\
\sem & 1.333 & 17.419 & 15.429 & 34.181 & 35 \\
\oth & 8.667 & 12.581 & 44.571 & 65.819 & 65 \\
\midrule
\multicolumn{6}{@{}l}{\emph{Converged (12 iterations, tolerance $10^{-10}$)}}\\
\sem & 1.381 & 17.717 & 15.902 & 35.000 & 35 \\
\oth & 8.619 & 12.283 & 44.098 & 65.000 & 65 \\
\bottomrule
\end{tabular}
\end{center}
Final semiconductor shares of imported division-26 inputs: $(0.138, 0.591, 0.265)$ versus the domestic
$(0.1, 0.5, 0.2)$. Trade data say this pair ships relatively more chips than domestic patterns imply, and RAS
raises every share while keeping their ranking.

\subsection{Stage 3: back to coefficients, eq.~(A.10)}
\begin{equation}
  a^{r,s}_{\sem,j} = \frac{\tZ^{r,s}_{\sem,j}}{Z^{r,s}_{\agg,j}}\,a^{r,s}_{\agg,j} = \frac{\tZ^{r,s}_{\sem,j}}{X^s_j},
  \qquad
  a^{r,s}_{\oth,j} = \frac{\tZ^{r,s}_{\oth,j}}{Z^{r,s}_{\agg,j}}\,a^{r,s}_{\agg,j} = \frac{\tZ^{r,s}_{\oth,j}}{X^s_j}.
  \tag{A.10}
\end{equation}
The first form reads: ``balanced share times original coefficient''. The second follows from
$a^{r,s}_{\agg,j} = Z^{r,s}_{\agg,j}/X^s_j$. Row additivity (Lemma~\ref{lem:rows}) holds by \eqref{eq:colj}.

%==========================================================================
\section{Order of Steps 3 and 4: the cell (\agg,\agg) of import blocks}\label{sec:order}
%==========================================================================

The import block $A^{r,s}$ contains the cell $a^{r,s}_{\agg,\agg}$ (division 26 of $r$ sold to division 26 of $s$),
which must become a $2\times 2$ block. It needs both a row and a column split.

\paragraph{Paper's order (column split first).} After (A.2), buyer columns \sem{} and \oth{} of $s$ exist, and the
RAS in Step 4 runs over $j$ including them. But Stage 1 needs \emph{flows}
$Z^{r,s}_{\agg,\sem} = a^{r,s}_{\agg,\sem}X^s_\sem$, i.e.\ it needs $X^s_\sem$, which is unknown until Step 5.
One would need a provisional $w^s$.

\paragraph{Recommended order (row split first).}
\begin{enumerate}[nosep]
  \item Step 4 on the \emph{original} columns $j=1,\dots,K$ (including the aggregate column \agg{} of $s$),
        using observed $Z^{r,s}_{\agg,j}$ and $\sigma^s_\agg = a^{s,s}_{\sem,\agg}/a^{s,s}_{\agg,\agg}$ for $j=\agg$.
        No unknown outputs are needed.
  \item Convert to coefficients with (A.10): this gives $a^{r,s}_{\sem,j}$, $a^{r,s}_{\oth,j}$ for all original $j$.
  \item Step 3 on all rows of $r$, now including \sem{} and \oth{}: for $k\in\{\sem,\oth\}$,
        $a^{r,s}_{k,l} = a^{r,s}_{k,\agg}\, a^{s,s}_{k,l}/a^{s,s}_{k,\agg}$, $l\in\{\sem,\oth\}$.
\end{enumerate}
The trade constraint \eqref{eq:rowS} is then matched over buyers including $s$'s whole division 26, which is then
distributed between \sem{} and \oth{} buyers by domestic ratios. Everything stays in coefficient space until
Step 5. \todo{confirm with authors/OECD if possible which order was used.}

%==========================================================================
\section{Step 5: final demand, output, flows and value added, eqs.~(A.11)--(A.13)}
%==========================================================================

\subsection{Final demand}
\paper{$Y^{r,s}_\sem = 0$ and $Y^{r,s}_\oth = Y^{r,s}_\agg$: semiconductors are only intermediate inputs.}
Apply to every final-demand category, including inventories (INVNT) and direct purchases abroad.

\subsection{Recomputing the table}
Stack the split coefficients into $A^\ast$ of dimension $N(K+1)$ and the final-demand row sums into $y^\ast$. Then
\begin{align}
  X^\ast &= (I-A^\ast)^{-1}y^\ast, \tag{A.11}\\
  Z^\ast &= A^\ast\hat X^\ast, \tag{A.12}\\
  V^\ast &= (I-A^{\ast\prime})\mathbf 1 . \tag{A.13}
\end{align}
(A.11) is the Leontief solution of $X = AX + y$. (A.12) inverts the definition $a_{ij}=Z_{ij}/X_j$.
(A.13) gives, for each column $j$, $1-\sum_i a_{ij}$: the share of output not spent on intermediates.

\begin{remark}[What $V$ contains]
In ICIO, $1-\sum_i a_{ij} = (T_j + W_j)/X_j$: it includes taxes less subsidies on intermediates as well as value
added. For unsplit industries use the original $T/W$ proportions. For \sem{} and \oth{} use the national-table
proportions (averaged as in Issue~\ref{iss:colavg} for non-reference economies).
\end{remark}

\subsection{Does the recomputed table reproduce the original?}\label{sec:checks}

\begin{proposition}
Suppose the split coefficients aggregate back to the original ones using the output weights $w^s_t$ that the
model itself produces, i.e.\ for every column Lemma~\ref{lem:cols} holds with $w^s = X^{\ast s}_\sem/X^{\ast s}_\agg$.
Then $X^\ast$ equals the original $X$ for every unsplit industry, and $X^{\ast s}_\sem + X^{\ast s}_\oth = X^s_\agg$.
\end{proposition}
\begin{proof}[Sketch]
Aggregating rows \sem, \oth{} and columns \sem, \oth{} (with weights $w$) of $X^\ast = A^\ast X^\ast + y^\ast$ yields
$\bar X = A\bar X + y$ for the aggregated vector $\bar X$; by uniqueness of the Leontief solution $\bar X = X$.
\end{proof}

The column splits aggregate with the \emph{base-year} weights $w_0$ (Step 1) or averaged weights $\bar w$
(Step 2), but the model produces its own $w_t = X^\ast_\sem/X^\ast_\agg$, driven by the row split and trade data.
Where $w_t\neq w_0$, outputs of \emph{other} industries move away from the published ICIO. The paper does not
discuss this. Required checks, per year:
\begin{enumerate}[nosep]
  \item $\max |X^\ast_i - X_i|/X_i$ over unsplit industries.
  \item $|X^{\ast s}_\sem + X^{\ast s}_\oth - X^s_\agg| / X^s_\agg$ per economy.
  \item Aggregated $Z^\ast$ vs.\ original $Z$, cell by cell.
  \item Non-negativity of $Z^\ast$, $V^\ast$ (negative $V$ signals intensities too high for a small $w$).
  \item Implied semiconductor exports $\sum_{s\neq r}\sum_j Z^{\ast r,s}_{\sem,j}$ vs.\ trade data.
  \item Implied world semiconductor VA and country shares vs.\ the paper's figures (e.g.\ five reference
        economies $\approx 75\%$ of world VA in 2018; CHN+KOR+TWN $\approx 60\%$).
\end{enumerate}

\paragraph{Optional flow-preserving reconciliation.} If the discrepancies matter for your use, keep original
$X$ for unsplit industries, set $X^s_\sem = \sum_{r,j}Z^{\ast s,r}_{\sem,j}$ (all sales are intermediate),
$X^s_\oth = X^s_\agg - X^s_\sem$, recompute column flows and rebalance the two new columns with RAS/GRAS against
the original column-26 flows. This departs from the paper and should be reported as such.

%==========================================================================
\section{Algorithm summary}
%==========================================================================

\begin{algorithm}[H]
\caption{Semiconductor split of ICIO (one year $t$)}
\begin{algorithmic}[1]
\Require ICIO $Z,Y,X$; trade $M^{r,s}_\sem$; ratios $\sigma^c,\rho^c,\beta^c,w^c_0$ for reference economies $c$
\State Build ratios for every economy: own ratios (reference) or averaged ratios (Issues~\ref{iss:avg}--\ref{iss:mys})
\State $A \gets Z\hat X^{-1}$
\For{each economy $r$} \Comment{Step 1--2: domestic block}
  \State split row 26 of $A^{r,r}$ with $\sigma^r$, column 26 with $\rho^r$, cell (26,26) with $\beta^r$
\EndFor
\For{each pair $(r,s)$, $r\neq s$} \Comment{Step 4: import rows}
  \State $Z^{(0)}_{\sem,j}\gets \sigma^s_j Z^{r,s}_{\agg,j}$, $Z^{(0)}_{\oth,j}\gets Z^{r,s}_{\agg,j}-Z^{(0)}_{\sem,j}$, $j=1..K$
  \State targets: rows $(M_\sem, \sum_j Z^{r,s}_{\agg,j}-M_\sem)$ (capped), columns $Z^{r,s}_{\agg,j}$
  \State $\tZ\gets\mathrm{RAS}(Z^{(0)},\text{targets})$; \quad $a^{r,s}_{k,j}\gets \tZ_{k,j}/X^s_j$
\EndFor
\For{each pair $(r,s)$, $r\neq s$} \Comment{Step 3: import columns}
  \State for every row $i$ of $r$ (incl.\ \sem, \oth): $a^{r,s}_{i,l}\gets a^{r,s}_{i,\agg}\,a^{s,s}_{i,l}/a^{s,s}_{i,\agg}$, $l\in\{\sem,\oth\}$
\EndFor
\State $y^\ast$: move all division-26 final demand to \oth; set \sem{} to 0
\State $X^\ast\gets(I-A^\ast)^{-1}y^\ast$; \quad $Z^\ast\gets A^\ast\hat X^\ast$; \quad $V^\ast\gets \mathbf 1 - A^{\ast\prime}\mathbf 1$
\State run checks of Section~\ref{sec:checks}
\end{algorithmic}
\end{algorithm}

%==========================================================================
\section{List of ambiguities and chosen resolutions}
%==========================================================================

\begin{center}
\begin{tabular}{@{}p{0.34\textwidth}p{0.6\textwidth}@{}}
\toprule
Issue & Resolution in this guide \\
\midrule
Typo in (A.2), \oth{} equation & First factor is $a^{r,s}_{i,\agg}$. \\
Iteration index $i$ in RAS & Renamed $k$. \\
\oth{} row target in RAS & $M_\oth = \sum_j Z^{r,s}_{\agg,j} - M_\sem$. \\
$M_\sem >$ intermediate div.-26 imports & Cap and log. \\
``Linear combination'' of coefficients & Average ratios, apply to own ICIO coefficients. \\
Averaging column ratios & Average flow shares $\phi$ and $w$, rebuild $\rho$. \\
Order of Steps 3/4 & Row split on original columns first. \\
Missing Malaysian table & See Issue~\ref{iss:mys}; report sensitivity. \\
$V$ mixes VA and TLS & Split using national-table proportions. \\
Recomputed $X$ differs from ICIO & Check; optional reconciliation. \\
Trade valuation (FOB/CIF) vs basic prices & Not adjusted in paper; document choice. \\
Footnote 10 ``described in 0'' & Broken cross-reference; refers to Step 2 rules. \\
\bottomrule
\end{tabular}
\end{center}

\section*{References}
\begin{itemize}[nosep,leftmargin=*]
  \item Bacharach, M. (1970), \emph{Biproportional Matrices and Input-Output Change}, Cambridge University Press.
  \item Haramboure, A., G. Lalanne, C. Schwellnus and J. Guilhoto (2023), ``Vulnerabilities in the semiconductor
        supply chain'', OECD Science, Technology and Industry Working Papers 2023/05, \url{https://doi.org/10.1787/6bed616f-en}.
  \item Martins Guilhoto, J., C. Webb and N. Yamano (2022), ``Guide to OECD TiVA Indicators, 2021 edition'',
        OECD STI Working Papers 2022/02, \url{https://doi.org/10.1787/58aa22b1-en}.
\end{itemize}

\end{document}
